\section{AHSME 1953}
        \begin{enumerate}
        	\item (\textit{AHSME 1953}) A boy buys oranges at $3$ for $10$ cents. He will sell them at $5$ for $20$ cents. In order to make a profit of $ $1.00$, he must sell: 


 $\textbf{(A)}\ 67 \text{ oranges} \qquad \textbf{(B)}\ 150 \text{ oranges} \qquad \textbf{(C)}\ 200\text{ oranges}\\  \textbf{(D)}\ \text{an infinite number of oranges}\qquad \textbf{(E)}\ \text{none of these}$

 \textbf{Answer: }B



\textbf{Solution}

The boy buys $3$ oranges for $10$ cents or $1$ orange for $\frac{10}{3}$ cents.  He sells them at $\frac{20}{5}=4$ cents each.  That means for every orange he sells, he makes a profit of $4-\frac{10}{3}=\frac{2}{3}$ cents.  

 To make a profit of $100$ cents, he needs to sell $\frac{100}{\frac{2}{3}}=\boxed{150}=\boxed{\text{B}}$




\textbf{Solution 2}

The boy buys $3$ oranges for $10$ cents. He sells them at $5$ for $20$ cents. So, he buys $15$ for $50$ cents and sells them at $15$ for $60$ cents, so he makes $10$ cents of profit on every $15$ oranges. To make $100$ cents of profit, he needs to sell $15 \cdot \frac{100}{10} = \boxed{150}$ oranges.

 


	\item (\textit{AHSME 1953}) A refrigerator is offered at sale at $250.00 less successive discounts of 20% and 15%. The sale price of the refrigerator is:


 $\textbf{(A) } \text{35\% less than 250.00} \qquad \textbf{(B) } \text{65\% of 250.00} \qquad \textbf{(C) } \text{77\% of 250.00} \qquad  \textbf{(D) } \text{68\% of 250.00} \qquad \textbf{(E) } \text{none of these}$

 \textbf{Answer: }D



\textbf{Solution}

The first discount takes off $20\%$ of the price, so the cost of the refrigerator is $0.8\cdot250$. The next discount takes off $15\%$, so the cost of the refrigerator is now $0.8\cdot0.85\cdot250=0.68\cdot250$. Thus, the sale price of the refrigerator is $\boxed{\textbf{(D)}\ \text{68\% of 250.00}}$.

 


	\item (\textit{AHSME 1953}) The factors of the expression $x^2+y^2$ are: 


 $\textbf{(A)}\ (x+y)(x-y) \qquad \textbf{(B)}\ (x+y)^2 \qquad \textbf{(C)}\ (x^{\frac{2}{3}}+y^{\frac{2}{3}})(x^{\frac{4}{3}}+y^{\frac{4}{3}})\\  \textbf{(D)}\ (x+iy)(x-iy)\qquad \textbf{(E)}\ \text{none of these}$

 \textbf{Answer: }D



\textbf{Solution}

Note that $x^2 + y^2 = x^2 - (-y^2) = x^2 - (i^2y^2) = x^2 - (iy)^2 = \boxed{\textbf{(D)}\ (x+iy)(x-iy)}$.

 


	\item (\textit{AHSME 1953}) The roots of $x(x^2+8x+16)(4-x)=0$ are: 


 $\textbf{(A)}\ 0 \qquad \textbf{(B)}\ 0,4 \qquad \textbf{(C)}\ 0,4,-4 \qquad \textbf{(D)}\ 0,4,-4,-4 \qquad \textbf{(E)}\ \text{none of these}$

 \textbf{Answer: }D



\textbf{Solution}

The equation, completely factored, is $x(x+4)^2(4-x)=0$.  Therefore, the roots, counting multiplicity, are therefore $\boxed{\textbf{(D)}\ 0,4,-4,-4 }$.

 


	\item (\textit{AHSME 1953}) If $\log_6 x=2.5$, the value of $x$ is: 


 $\textbf{(A)}\ 90 \qquad \textbf{(B)}\ 36 \qquad \textbf{(C)}\ 36\sqrt{6} \qquad \textbf{(D)}\ 0.5 \qquad \textbf{(E)}\ \text{none of these}$

 \textbf{Answer: }C



\textbf{Solution 1}

$\log_6 x=\frac{5}{2}\implies 2\log_6 x=5\implies \log_6 x^2=5\implies x^2=6^5\implies x=36\sqrt{6}$, $\fbox{C}$




\textbf{Solution 2}

$\log_6 x=\frac{5}{2}\implies 6^{\frac{5}{2}}=x \implies \sqrt{6^5}=x \implies 36\sqrt{6} = x$



	\item (\textit{AHSME 1953}) Charles has $5q + 1$ quarters and Richard has $q + 5$ quarters. The difference in their money in dimes is: 


 $\textbf{(A)}\ 10(q - 1) \qquad \textbf{(B)}\ \frac {2}{5}(4q - 4) \qquad \textbf{(C)}\ \frac {2}{5}(q - 1) \\  \textbf{(D)}\ \frac{5}{2}(q-1)\qquad \textbf{(E)}\ \text{none of these}$

 \textbf{Answer: }A



\textbf{Solution}

If Charles has $5q +1$ quarters and Richard has $q+5$ quarters, then the difference between them is  $4q-4$ quarters. Converting quarters to dimes, the difference between them is $\frac{5}{2}\cdot (4)(q-1)$, or $10(q-1)$ dimes, $\implies \fbox{A}$

 Note that $q$ does not stand for the number of quarters that they have.

 


	\item (\textit{AHSME 1953}) The fraction $\frac{\sqrt{a^2+x^2}-\frac{x^2-a^2}{\sqrt{a^2+x^2}}}{a^2+x^2}$ reduces to: 


 $\textbf{(A)}\ 0 \qquad \textbf{(B)}\ \frac{2a^2}{a^2+x^2} \qquad \textbf{(C)}\ \frac{2x^2}{(a^2+x^2)^{\frac{3}{2}}}\qquad \textbf{(D)}\ \frac{2a^2}{(a^2+x^2)^{\frac{3}{2}}}\qquad \textbf{(E)}\ \frac{2x^2}{a^2+x^2}$

 \textbf{Answer: }D



\textbf{Solution}

Multiplying the numerator and denominator by $\sqrt{a^2+x^2}$ results in 
 \[\frac{a^2+x^2-x^2+a^2}{(a^2+x^2)(\sqrt{a^2+x^2)}}=\frac{2a^2}{(a^2+x^2)(\sqrt{a^2+x^2)}}.\] Since $\sqrt{a^2+x^2}=(a^2+x^2)^{\frac{1}{2}}$, the denominator is $(a^2+x^2)^2\cdot (a^2+x^2)^{\frac{1}{2}} = (a^2+x^2)^{\frac{3}{2}}$ $\boxed{\textbf{(D) } \frac{2a^2}{(a^2+x^2)^{\frac{3}{2}}}}$.

 


	\item (\textit{AHSME 1953}) The value of $x$ at the intersection of $y=\frac{8}{x^2+4}$ and $x+y=2$ is: 


 $\textbf{(A)}\ -2+\sqrt{5} \qquad \textbf{(B)}\ -2-\sqrt{5} \qquad \textbf{(C)}\ 0 \qquad \textbf{(D)}\ 2 \qquad \textbf{(E)}\ \text{none of these}$

 \textbf{Answer: }C



\textbf{Solution}

Isolating $y$ in the second equation yields $y = 2-x$.  Substituting for $y$ in the first equation yields $2-x=\frac{8}{x^2+4}$, which can be rearranged to form $x(x^2-2x+4)=0$.  The only value of $x$ satisfying this equation is $x = \boxed{\textbf{(C)}\ 0}$, since $x^2-2x+4$ has no real roots.

 


	\item (\textit{AHSME 1953}) The number of ounces of water needed to reduce $9$ ounces of shaving lotion containing $50$ \% alcohol to a lotion containing $30$ \% alcohol is: 


 $\textbf{(A)}\ 3 \qquad \textbf{(B)}\ 4 \qquad \textbf{(C)}\ 5 \qquad \textbf{(D)}\ 6 \qquad \textbf{(E)}\ 7$

 \textbf{Answer: }D



\textbf{Solution 1}

Say we add $N$ ounces of water to the shaving lotion. Since half of a $9$ ounce bottle of shaving lotion is alcohol, we know that we have $\frac{9}{2}$ ounces of alcohol. We want $\frac{9}{2}=0.3(9+N)$ (because we want the amount of alcohol, $\frac{9}{2}$, to be $30\%$, or $0.3$, of the total amount of shaving lotion, $9+N$). Solving, we find that 
 \[9=0.6(9+N)\implies9=5.4+0.6N\implies3.6=0.6N\implies6=N.\] So, the total amount of water we need to add is $\boxed{\textbf{(D) } 6}$.

 



\textbf{Solution 2}

The concentration of alcohol after adding $n$ ounces of water is $\frac{4.5}{9+n}$. 

 To get a solution of 30\% alcohol, we solve $\frac{4.5}{9+n}=\frac{3}{10}$
$45=27+3n$
$18=3n$
$6=n \implies \boxed{\textbf{(D) } 6}$



	\item (\textit{AHSME 1953}) The number of revolutions of a wheel, with fixed center and with an outside diameter of $6$ feet, required to cause a point on the rim to go one mile is: 


 $\textbf{(A)}\ 880 \qquad \textbf{(B)}\ \frac{440}{\pi} \qquad \textbf{(C)}\ \frac{880}{\pi} \qquad \textbf{(D)}\ 440\pi\qquad \textbf{(E)}\ \text{none of these}$

 \textbf{Answer: }C



\textbf{Solution}

We know that the radius of the wheel is $3$ feet, so the total circumference of the wheel is $6\pi$ feet. We also know that one mile is equivalent to $5280$ feet. It takes $\frac{5280}{6\pi}$ revolutions for any one point on the wheel to travel a mile. Simplifying, we find that the answer is $\boxed{\textbf{(C) } \frac{880}{\pi}}$.

 


 


	\item (\textit{AHSME 1953}) A running track is the ring formed by two concentric circles. It is $10$ feet wide. The circumference of the two circles differ by about: 


 $\textbf{(A)}\ 10\text{ feet} \qquad \textbf{(B)}\ 30\text{ feet} \qquad \textbf{(C)}\ 60\text{ feet} \qquad \textbf{(D)}\ 100\text{ feet}\\ \textbf{(E)}\ \text{none of these}$

 \textbf{Answer: }C



\textbf{Solution 1}

Let $r$ be the radius of the inner circle.  Then the outer circle has radius $r + 10$, so the difference between their circumferences is $2\pi(r+10) - 2\pi = 2\pi \cdot 10 = 20\pi \approx \boxed{\textbf{(C)}\ 60\text{ feet} }$




\textbf{Solution 2}

Since the track is 10 feet wide, the diameter of the outer circle will be 20 feet more than the inner circle. Since the circumference of a circle is directly proportional to its diameter, the difference in the circles' diameters is simply $20\pi$ feet. Using $\pi \approx 3$, the answer is $\fbox{C}$.

 


	\item (\textit{AHSME 1953}) The diameters of two circles are $8$ inches and $12$ inches respectively. The ratio of the area of the smaller to the area of the larger circle is: 


 $\textbf{(A)}\ \frac{2}{3} \qquad \textbf{(B)}\ \frac{4}{9} \qquad \textbf{(C)}\ \frac{9}{4} \qquad \textbf{(D)}\ \frac{1}{2}\qquad \textbf{(E)}\ \text{none of these}$

 \textbf{Answer: }B



\textbf{Solution 1}

The smaller circle has area $\pi\cdot4^2 = 16\pi$ while the larger circle has area $\pi\cdot6^2=36\pi$. So the ratio of the area of the smaller circle to the area of the larger circle is $\frac{16\pi}{36\pi}=\boxed{\textbf{(B) }\frac{4}{9}}$.

 



\textbf{Solution 2}

Recall that the ratio between the areas of any two similar figures must equal the square of the ratio between their corresponding dimensions.  Thus the ratio between the areas of these two circles is $\left(\frac{8}{12}\right)^2 = \left(\frac{2}{3}\right)^2= \boxed{\textbf{(B) }\frac{4}{9}}$.

 


	\item (\textit{AHSME 1953}) A triangle and a trapezoid are equal in area. They also have the same altitude. If the base of the triangle is 18 inches, the median of the trapezoid is: 


 $\textbf{(A)}\ 36\text{ inches} \qquad \textbf{(B)}\ 9\text{ inches} \qquad \textbf{(C)}\ 18\text{ inches}\\  \textbf{(D)}\ \text{not obtainable from these data}\qquad \textbf{(E)}\ \text{none of these}$

 \textbf{Answer: }B



\textbf{Solution}

We know that the area of a trapezoid is $mh$ where $m$ is the median and the area of a triangle is $\frac{bh}{2}$. We know that $b=18$ and the height of the trapezoid is congruent to the height of the triangle. Thus, we get $\frac{18h}{2}=9h=mh$, so $m=\textbf{(B) }9\text{ inches}$.

 


 


	\item (\textit{AHSME 1953}) Given the larger of two circles with center $P$ and radius $p$ and the smaller with center $Q$ and radius $q$. Draw $PQ$. Which of the following statements is false? 


 $\textbf{(A)}\ p-q\text{ can be equal to }\overline{PQ}\\  \textbf{(B)}\ p+q\text{ can be equal to }\overline{PQ}\\  \textbf{(C)}\ p+q\text{ can be less than }\overline{PQ}\\  \textbf{(D)}\ p-q\text{ can be less than }\overline{PQ}\\ \textbf{(E)}\ \text{none of these}$

 \textbf{Answer: }E



\textbf{Solution 1}

Let $r$ be any positive real number.  Draw a segment of length $r$, and label the ends $P$ and $Q$.  Draw a circle of radius $p$ centered at $P$ and a circle of radius $q$ centered at $Q$.  No matter what $r$ is, this is possible to do.  

 To show that choice A can happen, use $r = p-q$.  For choice B, use $r = p+q$.  For choice C, use $r = p+q+1$.  For choice D, use $r = p-q+1$. So $\boxed{\textbf{(E)}\ \text{none of these} }$ is correct.

 



\textbf{Solution 2}

We will test each option to see if it can be true or not. Links to diagrams are provided.
 \[\textbf{(A)}\ p-q\text{ can be equal to }\overline{PQ}\]Let circle $Q$ be inside circle $P$ and tangent to circle $P$, and the point of tangency be $R$. $PR = p$, and $QR = q$, so $PR - QR = PQ = p-q.$
                    \begin{figure}[H]
                        \centering
                        \includegraphics[width=0.5\textwidth]{Solution/AHSME_I_1953_14_Solution_Figure_1.png}
                    \end{figure}
                    
 \[\textbf{(B)}\ p+q\text{ can be equal to }\overline{PQ}\]Let circle $Q$ be outside circle $P$ and tangent to circle $P$, and the point of tangency be $R$. $PR = p$, and $QR = q$, so $PR + QR = PQ = p+q.$
                    \begin{figure}[H]
                        \centering
                        \includegraphics[width=0.5\textwidth]{Solution/AHSME_I_1953_14_Solution_Figure_2.png}
                    \end{figure}
                    
 \[\textbf{(C)}\ p+q\text{ can be less than }\overline{PQ}\]Let circle $Q$ be outside circle $P$ and not tangent to circle $P$, and the intersection of $\overline{PQ}$ with the circles be $R$ and $S$ respectively. $PR = p$ and $QS = q$, and $PR + QS < PQ$, so $p+q < PQ.$
                    \begin{figure}[H]
                        \centering
                        \includegraphics[width=0.5\textwidth]{Solution/AHSME_I_1953_14_Solution_Figure_3.png}
                    \end{figure}
                    
 \[\textbf{(D)}\ p-q\text{ can be less than }\overline{PQ}\]In the situation for choice C, $p+q < PQ$.  Since $p-q < p+q$, $p-q < PQ$.  So this is possible.

 Since options A, B, C, and D can be true, the answer must be $\boxed{E}$.

 


	\item (\textit{AHSME 1953}) A circular piece of metal of maximum size is cut out of a square piece and then a square piece of maximum size is cut out of the circular piece. The total amount of metal wasted is: 


 $\textbf{(A)}\ \frac{1}{4} \text{ the area of the original square}\\  \textbf{(B)}\ \frac{1}{2}\text{ the area of the original square}\\ \textbf{(C)}\ \frac{1}{2}\text{ the area of the circular piece}\\ \textbf{(D)}\ \frac{1}{4}\text{ the area of the circular piece}\\ \textbf{(E)}\ \text{none of these}$

 \textbf{Answer: }B



\textbf{Solution}

The maximum diameter of the circular piece is the same as the side length of the square piece, so the circle is tangent to the square on all four sides. The maximum size a square piece that you can cut from the circle now has 4 edges that are the same as the square created by connecting the four midpoints of the original square. If the side length of the original square is $s$, then the sides of the final square have length $\sqrt{\frac{s^2}{2}}$, so its area is $\frac{s^2}{2}$, $\frac{1}{2}$ the area of the original square. $\textbf{(B)}$



	\item (\textit{AHSME 1953}) Adams plans a profit of $10$ \% on the selling price of an article and his expenses are $15$ \% of sales. The rate of markup on an article that sells for $ $5.00$ is: 


 $\textbf{(A)}\ 20\% \qquad \textbf{(B)}\ 25\% \qquad \textbf{(C)}\ 30\% \qquad \textbf{(D)}\ 33\frac {1}{3}\% \qquad \textbf{(E)}\ 35\%$

 \textbf{Answer: }D



\textbf{Solution}

Let the cost of the product be C, and the selling price be S. Since Adams wants a 10\% profit on the sale, and his expenses are 15\% of the sale, we add $0.10S$ and $0.15S$ to the cost to get the selling price.Then, $S=C+0.15S+0.10S \Rightarrow 0.75S=C \Rightarrow S=\frac{4}{3}C \Rightarrow \fbox{D}$.

 


	\item (\textit{AHSME 1953}) A man has part of $ $4500$ invested at $4$ % and the rest at $6$ %. If his annual return on each investment is the same, the average rate of interest which he realizes of the $4500 is: 


 $\textbf{(A)}\ 5\% \qquad \textbf{(B)}\ 4.8\% \qquad \textbf{(C)}\ 5.2\% \qquad \textbf{(D)}\ 4.6\% \qquad \textbf{(E)}\ \text{none of these}$

 \textbf{Answer: }B



\textbf{Solution}

Let $x$ be the amount that was invested at $6\%$, so $4500-x$ was invested at $4\%$.  Then $0.04(4500-x) = 0.06x$.  Solving this equation yields $x = 1800$.  Therefore, the total return on the investments is $0.04\cdot2700+0.06\cdot1800=216$, for an average interest rate of $\frac{216}{4500} = 0.048 = \boxed{\textbf{(B)}\ 4.8\%}$.

 


	\item (\textit{AHSME 1953}) One of the factors of $x^4+4$ is: 


 $\textbf{(A)}\ x^2+2 \qquad \textbf{(B)}\ x+1 \qquad \textbf{(C)}\ x^2-2x+2 \qquad \textbf{(D)}\ x^2-4\\  \textbf{(E)}\ \text{none of these}$

 \textbf{Answer: }C



\textbf{Solution 1}

Note that $x^4 + 4 = x^4 - 4x^2 + 4 + 4x^2 = (x^2 + 2)^2 - (2x)^2 = (x^2 - 2x + 2)(x^2 + 2x + 2)$.  So $\boxed{\textbf{(C)}\ x^2-2x+2}$ is correct.

 



\textbf{Solution 2}

Recall the Sophie Germain Identity $a^4 + 4b^4 = (a^2 - 2ab + 2b^2)(a^2 + 2ab + 2b^2)$.

 So $x^4 + 4 = (x^2-2x+2)(x^2+2x+2)$, and $\boxed{\textbf{(C)}\ x^2-2x+2}$ is correct.

 



\textbf{Solution 3 (Answer Choices)}

Choice (A) is incorrect, as the roots of $x^2 + 2$ are $\pm i\sqrt{2}$, neither of which is a root of $x^4 + 4$.

 Choice (B) is incorrect, as the root of $x + 1$ is $-1$, which is not a root of $x^4 + 4$.

 Choice (D) is incorrect, as the roots of $x^2 - 4$ are $\pm 2$, neither of which is a root of $x^4 + 4$.

 
Choice (C) is correct, as the roots of $x^2 - 2x + 2$ are $1 \pm i$, both of which are roots of $x^4 + 4$. 

 


	\item (\textit{AHSME 1953}) In the expression $xy^2$, the values of $x$ and $y$ are each decreased $25$ \%; the value of the expression is: 


 $\textbf{(A)}\ \text{decreased } 50\% \qquad \textbf{(B)}\ \text{decreased }75\%\\  \textbf{(C)}\ \text{decreased }\frac{37}{64}\text{ of its value}\qquad \textbf{(D)}\ \text{decreased }\frac{27}{64}\text{ of its value}\\ \textbf{(E)}\ \text{none of these}$

 \textbf{Answer: }C



\textbf{Solution}

If $x$ and $y$ are decreased by $25\%$, $x$ and $y$ become $\frac{3}{4}x$ and $\frac{3}{4}y$ respectively.  So $xy^2$ becomes $\left(\frac{3}{4}x\right)\left(\frac{3}{4}y\right)^2 = \frac{27}{64}xy^2$, so it decreases by $xy^2 - \frac{27}{64}xy^2 = \frac{37}{64}xy^2$.  Therefore, $xy^2$ is $\boxed{\textbf{(C)}\ \text{decreased }\frac{37}{64}\text{ of its value}}$.

 


	\item (\textit{AHSME 1953}) If $y=x+\frac{1}{x}$, then $x^4+x^3-4x^2+x+1=0$ becomes: 


 $\textbf{(A)}\ x^2(y^2+y-2)=0 \qquad \textbf{(B)}\ x^2(y^2+y-3)=0\\  \textbf{(C)}\ x^2(y^2+y-4)=0 \qquad \textbf{(D)}\ x^2(y^2+y-6)=0\\ \textbf{(E)}\ \text{none of these}$

 \textbf{Answer: }D



\textbf{Solution 1}

$x^4+x^3-4x^2+x+1 = x^2\left(x^2 + x - 4 + \frac{1}{x} + \frac{1}{x^2}\right) = x^2\left(x^2 + 2 + \frac{1}{x^2} + x + \frac{1}{x} - 6\right) = x^2\left(y^2+y-6\right)$.

 



\textbf{Solution 2}

We multiply each of the answers to get: $x^2(y^2)+x^2(y)+nx^2$, where $n$ is either $-2,-3,-4,$ or $-6$. Looking at the first term, we have to square $y$, or $x+\frac{1}{x}$, doing so, we get the equation $x^2+\frac{1}{x^2}+2$. Multiplying that by $x^2$, we get $x^4+2x^2+1$. Multiplying $y$ by $x^2$, we get the expression $x^3+x$. Adding these two equations together, we get $x^4+x^3+2x^2+x+1+nx=0$. To get the term $-4x^2$, which was in the original equation, $n$ must be $-6$, giving an answer of $\boxed{D}$



	\item (\textit{AHSME 1953}) If $\log_{10} (x^2-3x+6)=1$, the value of $x$ is: 


 $\textbf{(A)}\ 10\text{ or }2 \qquad \textbf{(B)}\ 4\text{ or }-2 \qquad \textbf{(C)}\ 3\text{ or }-1 \qquad \textbf{(D)}\ 4\text{ or }-1\\ \textbf{(E)}\ \text{none of these}$

 \textbf{Answer: }D



\textbf{Solution}

We know that $x^2-3x+6=10^1$, after expanding the logarithm. Taking the 10 to the other side, $x^2-3x-4=0$. Factoring, we get $(x-4)(x+1)=0$, so $x=\boxed{4,-1\Rightarrow \text{(D)}}$.

 


	\item (\textit{AHSME 1953}) The logarithm of $27\sqrt[4]{9}\sqrt[3]{9}$ to the base $3$ is: 


 $\textbf{(A)}\ 8\frac{1}{2} \qquad \textbf{(B)}\ 4\frac{1}{6} \qquad \textbf{(C)}\ 5 \qquad \textbf{(D)}\ 3 \qquad \textbf{(E)}\ \text{none of these}$

 \textbf{Answer: }B



\textbf{Solution}

Rewriting as powers of $3$ yields $27\sqrt[4]{9}\sqrt[3]{9} = 3^3\cdot 3^\frac{1}{2}\cdot 3^\frac{2}{3} = 3^{3+\frac{1}{2}+\frac{2}{3}} = 3^\frac{25}{6}$.So $\log_3\left({27\sqrt[4]{9}\sqrt[3]{9}}\right) = \log_3{3^\frac{25}{6}} = \frac{25}{6} = \boxed{\text{(B)}\ 4\frac{1}{6}}$.

 


	\item (\textit{AHSME 1953}) The equation $\sqrt {x + 10} - \frac {6}{\sqrt {x + 10}} = 5$ has: 


 $\textbf{(A)}\ \text{an extraneous root between } - 5\text{ and } - 1 \\  \textbf{(B)}\ \text{an extraneous root between }-10\text{ and }-6\\ \textbf{(C)}\ \text{a true root between }20\text{ and }25\qquad \textbf{(D)}\ \text{two true roots}\\ \textbf{(E)}\ \text{two extraneous roots}$

 \textbf{Answer: }B



\textbf{Solution 1}

Let $u = \sqrt{x+10}$.  Then $u - \frac{6}{u} = 5$, that is, $u^2 - 5u - 6 = (u-6)(u+1) = 0$.  So $\sqrt{x+10} = 6$ or $\sqrt{x + 10} = -1$.  

 The equation $\sqrt{x+10} = 6$ has a real solution of $x = 26$, while $\sqrt{x+10} = -1$ has an extraneous solution of $x = -9$.  Therefore, the correct answer is $\boxed{\textbf{(B)}\ \text{an extraneous root between }-10\text{ and }-6}$.

 



\textbf{Solution 2}

We can multiply both sides by $\sqrt{x+10}$ to get $x+4=5\sqrt{x+10}$. We can now square both sides to get $x^2+8x+16=25x+250$, which yields $x^2-17x-234=0$. We can factor the quadratic as $(x+9)(x-26)=0$, giving us roots of $-9$ and $26$. Plugging in these values, we find that $-9$ is an extraneous root and $26$ is a true root. This gives us the final answer of $\boxed{\textbf{(B)}\ \text{an extraneous root between }-10\text{ and }-6}$.

 


	\item (\textit{AHSME 1953}) If $a,b,c$ are positive integers less than $10$, then $(10a + b)(10a + c) = 100a(a + 1) + bc$ if: 


 $\textbf{(A)}\ b + c = 10 \qquad \textbf{(B)}\ b = c \qquad \textbf{(C)}\ a + b = 10 \qquad \textbf{(D)}\ a = b \\  \textbf{(E)}\ a+b+c = 10$

 \textbf{Answer: }A



\textbf{Solution}

Multiply out the LHS to get $100a^2+10ac+10ab+bc=100a(a+1)+bc$. Subtract $bc$ and factor to get $10a(10a+b+c)=10a(10a+10)$. Divide both sides by $10a$ and then subtract $10a$ to get $b+c=10$, giving an answer of $\boxed{A}$.

 


	\item (\textit{AHSME 1953}) In a geometric progression whose terms are positive, any term is equal to the sum of the next two following terms. then the common ratio is: 


 $\textbf{(A)}\ 1 \qquad \textbf{(B)}\ \text{about }\frac{\sqrt{5}}{2} \qquad \textbf{(C)}\ \frac{\sqrt{5}-1}{2}\qquad \textbf{(D)}\ \frac{1-\sqrt{5}}{2}\qquad \textbf{(E)}\ \frac{2}{\sqrt{5}}$

 \textbf{Answer: }C



\textbf{Solution}

Let $a$ be the first term and $r$ be the common ration of this sequence. Then $a=ar+ar^2$, and therefore $r^2+r-1=0$. By the quadratic formula, the solutions to this equation are $r = \frac{-1\pm\sqrt{5}}{2}$.  Since $r$ must be positive, $r = \boxed{\textbf{(C)}\ \frac{\sqrt{5}-1}{2}}$.

 


	\item (\textit{AHSME 1953}) The base of a triangle is $15$ inches. Two lines are drawn parallel to the base, terminating in the other two sides, and dividing the triangle into three equal areas. The length of the parallel closer to the base is: 


 $\textbf{(A)}\ 5\sqrt{6}\text{ inches} \qquad \textbf{(B)}\ 10\text{ inches} \qquad \textbf{(C)}\ 4\sqrt{3}\text{ inches}\qquad \textbf{(D)}\ 7.5\text{ inches}\\  \textbf{(E)}\ \text{none of these}$

 \textbf{Answer: }A



\textbf{Solution}

Let the triangle be $\triangle ABC$ where $BC$ is the base. Then let the parallels be $MN$ and $PQ$, where $PQ$ is closer to our base $BC$. 

 It's obvious that $\triangle APQ \sim \triangle ABC$, where $|\triangle APQ|:|\triangle ABC|=2:3$, so $\frac{PQ}{BC}=\sqrt{\frac{2}{3}}$. Since we know $BC=15$, $PQ=\frac{15\sqrt{2}}{\sqrt{3}}=5\sqrt{6}$

 Hence our answer is $\textbf{(A)}$



	\item (\textit{AHSME 1953}) The radius of the first circle is $1$ inch, that of the second $\frac{1}{2}$ inch, that of the third $\frac{1}{4}$ inch and so on indefinitely. 
The sum of the areas of the circles is: 


 $\textbf{(A)}\ \frac{3\pi}{4} \qquad \textbf{(B)}\ 1.3\pi \qquad \textbf{(C)}\ 2\pi \qquad \textbf{(D)}\ \frac{4\pi}{3}\qquad \textbf{(E)}\ \text{none of these}$

 \textbf{Answer: }D



\textbf{Solution}

Note the areas of these circles is $1\pi$, $\frac{\pi}{4}$, $\frac{\pi}{16}, \dots$. The sum of these areas will thus be $\pi\left(1+\frac{1}{4}+\frac{1}{16}+\dots\right)$. We use the formula for an infinite geometric series to get the sum of the areas will be $\pi\left(\frac{1}{1-\frac{1}{4}}\right)$, or $\frac{4\pi}{3}$.  $\boxed{D}$



	\item (\textit{AHSME 1953}) In $\triangle ABC$, sides $a,b$ and $c$ are opposite $\angle{A},\angle{B}$ and $\angle{C}$ respectively. $AD$ bisects $\angle{A}$ and meets $BC$ at $D$. 
Then if $x = \overline{CD}$ and $y = \overline{BD}$ the correct proportion is: 


 $\textbf{(A)}\ \frac {x}{a} = \frac {a}{b + c} \qquad \textbf{(B)}\ \frac {x}{b} = \frac {a}{a + c} \qquad \textbf{(C)}\ \frac{y}{c}=\frac{c}{b+c}\\ \textbf{(D)}\ \frac{y}{c}=\frac{a}{b+c}\qquad \textbf{(E)}\ \frac{x}{y}=\frac{c}{b}$

 \textbf{Answer: }D



\textbf{Solution}


                    \begin{figure}[H]
                        \centering
                        \includegraphics[width=0.5\textwidth]{Solution/AHSME_I_1953_28_Solution_Figure_1.png}
                    \end{figure}
                    By the \href{https://artofproblemsolving.com/wiki/index.php?title=Angle\_Bisector\_Theorem}{Angle Bisector Theorem}, $\frac{x}{b}=\frac{y}{c}$. Because of this, $\frac{x}{b}$ must equal $\frac{ky}{kc}$, where $k$ is some real number. Therefore, $\frac{x+y}{b+c}=\frac{y+ky}{c+kc}$. Factoring $\frac{k+1}{k+1}$ out, we get $\frac{y}{c}=\frac{x+y}{b+c}$. However, $x+y=a$, so the answer is $\textbf{(D)}\ \frac{y}{c}=\frac{a}{b+c}$.

 


	\item (\textit{AHSME 1953}) The number of significant digits in the measurement of the side of a square whose computed area is $1.1025$ square inches to 
the nearest ten-thousandth of a square inch is: 


 $\textbf{(A)}\ 2 \qquad \textbf{(B)}\ 3 \qquad \textbf{(C)}\ 4 \qquad \textbf{(D)}\ 5 \qquad \textbf{(E)}\ 1$

 \textbf{Answer: }D



\textbf{Solution 1}

The area of the square must be between $1.10245$ and $1.10255$.  

 Let $s$ be the side length of the square, and let $x = s - 1.05$.  Then the area of the square is $(1.05 + x)^2 = 1.1025 + 2.1x + x^2$.

 If $x = -0.00005$, then the area is $1.1025 - 2.1 \cdot 0.00005 + 0.00005^2 = 1.1025 - 2.09995 \cdot 0.00005 < 1.1024$. If $x = 0.00005$, then the area is $1.1025 + 2.1 \cdot 0.00005 + 0.00005^2 > 1.1026$.

 Therefore, $s=1.04995$ is too low and $x = 1.05005$ is too high for $s^2$ to be correct to $4$ decimal places, so $1.04995 < s < 1.05005$, so $s = 1.0500$ must be accurate to $\boxed{\textbf{(D)}\ 5}$ significant digits.

 However, $s$ may not be accurately determined to $6$ significant digits, as $1.05002^2 = 1.1025 + 2.1 \cdot 0.00002 + 0.00002^2 = 1.1025 + 2.10002 \cdot 0.00002 < 1.10255$, so $s$ could be anywhere between $1.05000$ and $1.05002$.

 



\textbf{Solution 2}

Note that $\sqrt{1.1025} = 1.05$.  Let $s$ be the side length of the square and let $x = s - 1.05$.  Then the area of the square is $(1.05 + x)^2 = 1.1025 + 2.1x + x^2 \approx 1.1025 + 2.1x$ ($x^2$ will be negligible since $x$ is already small).    

 Since the area of the square is accurate to $4$ decimal places, $2.1|x| < 0.00005$ and thus $|x| < 0.00005$.  That is, $1.04995 < s < 1.05005$, so the side length may be determined accurately to $\boxed{\textbf{(D)}\ 5}$ significant digits.

 


	\item (\textit{AHSME 1953}) A house worth $ $9000$ is sold by Mr. A to Mr. B at a $10$ % loss. Mr. B sells the house back to Mr. A at a $10$ % gain. 
The result of the two transactions is: 


 $\textbf{(A)}\ \text{Mr. A breaks even} \qquad \textbf{(B)}\ \text{Mr. B gains }\$900 \qquad \textbf{(C)}\ \text{Mr. A loses }\$900\\ \textbf{(D)}\ \text{Mr. A loses }\$810\qquad \textbf{(E)}\ \text{Mr. B gains }\$1710$

 \textbf{Answer: }D



\textbf{Solution}

When Mr.A sells the house at a $10\%$ loss, he sells it for $9000(1 - .1) = 8100$. When Mr.B sells the house back to Mr. A at a $10\%$ gain he sells it for $8100(1 + .1) = 8910$. Therefore Mr. A has lost $8100-8910 = 810$ dollars, so the answer is $\boxed{D}$.

 


	\item (\textit{AHSME 1953}) The rails on a railroad are $30$ feet long. As the train passes over the point where the rails are joined, there is an audible click. 
The speed of the train in miles per hour is approximately the number of clicks heard in: 


 $\textbf{(A)}\ 20\text{ seconds} \qquad \textbf{(B)}\ 2\text{ minutes} \qquad \textbf{(C)}\ 1\frac{1}{2}\text{ minutes}\qquad \textbf{(D)}\ 5\text{ minutes}\\ \textbf{(E)}\ \text{none of these}$

 \textbf{Answer: }A



\textbf{Solution}

Let $x$ be the speed of the train in miles per hour. By dimensional analysis, the speed in clicks per minute is:
 \[\left(\frac{x\text{ mi}}{1\text{ hr}}\right)\left(\frac{1\text{ hr}}{60\text{ min}}\right)\left(\frac{5280\text{ ft}}{1\text{ mi}}\right)\left(\frac{1\text{ click}}{30\text{ ft}}\right)=\frac{5280x\text{ click}}{1800\text{ min}} = \frac{44}{15}x \text{ click}/\text{min} \approx 3x\text{ click}/\text{min}\]

 Therefore, approximately $3x$ clicks are heard per minute. So approximately $x$ clicks are heard per $\frac{1}{3}$ minute, that is, approximately $x$ clicks are heard per $\boxed{\textbf{(A)}\ 20\text{ seconds}}$.

 


	\item (\textit{AHSME 1953}) Each angle of a rectangle is trisected. The intersections of the pairs of trisectors adjacent to the same side always form: 


 $\textbf{(A)}\ \text{a square} \qquad \textbf{(B)}\ \text{a rectangle} \qquad \textbf{(C)}\ \text{a parallelogram with unequal sides}\\ \textbf{(D)}\ \text{a rhombus}\qquad \textbf{(E)}\ \text{a quadrilateral with no special properties}$

 \textbf{Answer: }D



\textbf{Solution}


                    \begin{figure}[H]
                        \centering
                        \includegraphics[width=0.5\textwidth]{Solution/AHSME_I_1953_32_Solution_Figure_1.png}
                    \end{figure}
                    

 Let $A$, $B$, $C$, and $D$ be the vertices of the rectangle, and let $E$, $F$, $G$ and $H$ be the intersections of the angle trisectors.  

 Since $ABCD$ is a rectangle, $AB = CD$ and $AD = BC$.  Since $\angle EAB = \angle EBA = \angle GCD = \angle GDC = 30^{\circ}$, $\triangle AEB$ and $\triangle CGD$ are isosceles.  Also, by ASA, $\triangle AEB \cong \triangle CGD$, and thus $AE = EB = CG = GD$.

 Also, $\angle FBC = \angle FCB = \angle HAD = \angle HDA = 30^{\circ}$, so $\triangle BFC$ and $\triangle AHD$ are also isosceles. By ASA, $\triangle BFC \cong \triangle AHD$, and thus $BF = FC = AH = HD$.

 Since $\angle HAE = \angle FBE = \angle FCG = \angle HDG = 30^{\circ}$, by SAS, $\triangle HAE \cong \triangle FBE \cong \triangle FCG = \triangle HDG$, and thus $HE = FE = FG = HG$, so $EFGH$ is $\boxed{\textbf{(D)}\ \text{a rhombus}}$. 

 \textbf{Remark} If $\frac{BC}{AB} = \sqrt{3}$, then $F = H$ and the rhombus formed is degenerate.

 


	\item (\textit{AHSME 1953}) The perimeter of an isosceles right triangle is $2p$. Its area is: 


 $\textbf{(A)}\ (2+\sqrt{2})p \qquad \textbf{(B)}\ (2-\sqrt{2})p \qquad \textbf{(C)}\ (3-2\sqrt{2})p^2\\  \textbf{(D)}\ (1-2\sqrt{2})p^2\qquad \textbf{(E)}\ (3+2\sqrt{2})p^2$

 \textbf{Answer: }C



\textbf{Solution}

Given leg length $x$, we can write the perimeter of this triangle to be $2x+x\sqrt{2}=2p$. Thus, $x(2+\sqrt{2})=2p$. Divide to get $x=\frac{2p}{2+\sqrt{2}}$. Multiply by the conjugate and simplify to get $x=p(2-\sqrt{2})$. Square and divide by two to get the area of the triangle, or $\frac{p^2*(2-\sqrt{2}^2)}{2}$, or $p^2(3-2\sqrt{2})$. $\boxed{C}$



	\item (\textit{AHSME 1953}) If one side of a triangle is $12$ inches and the opposite angle is $30^{\circ}$, then the diameter of the circumscribed circle is: 


 $\textbf{(A)}\ 18\text{ inches} \qquad \textbf{(B)}\ 30\text{ inches} \qquad \textbf{(C)}\ 24\text{ inches} \qquad \textbf{(D)}\ 20\text{ inches}\\ \textbf{(E)}\ \text{none of these}$

 \textbf{Answer: }C



\textbf{Solution}

By the Extended Law of Sines, the diameter, or twice the circumradius $R$, is given by
 \[2R=\frac{12\text{ inches}}{\sin30^\circ}=\boxed{\textbf{(C)}\ 24\text{ inches}}.\]



	\item (\textit{AHSME 1953}) If $f(x)=\frac{x(x-1)}{2}$, then $f(x+2)$ equals: 


 $\textbf{(A)}\ f(x)+f(2) \qquad \textbf{(B)}\ (x+2)f(x) \qquad \textbf{(C)}\ x(x+2)f(x) \qquad \textbf{(D)}\ \frac{xf(x)}{x+2}\\ \textbf{(E)}\ \frac{(x+2)f(x+1)}{x}$

 \textbf{Answer: }E



\textbf{Solution}

First simplify $f(x+2)$:
 \[f(x+2)=\frac{(x+2)(x+1)}{2}.\]Then simplify the answers and see which one matches.

 $f(x)+f(2)=\frac{x(x-1)}{2} + 1,$
$(x+2)f(x)=\frac{(x+2)x(x-1)}{2},$
$x(x+2)f(x)=\frac{(x+2)x^2(x-1)}{2},$
$\frac{xf(x)}{x+2}=\frac{x^2(x-1)}{2(x+2)},$
$\frac{(x+2)f(x+1)}{x}=\frac{(x+2)(x+1)}{2}.$

 Since $f(x+2)=\frac{(x+2)f(x+1)}{x}=\frac{(x+2)(x+1)}{2}$, the answer is $\boxed{\textbf{(E)}\ \frac{(x+2)f(x+1)}{x}}.$



	\item (\textit{AHSME 1953}) Determine $m$ so that $4x^2-6x+m$ is divisible by $x-3$. The obtained value, $m$, is an exact divisor of: 


 $\textbf{(A)}\ 12 \qquad \textbf{(B)}\ 20 \qquad \textbf{(C)}\ 36 \qquad \textbf{(D)}\ 48 \qquad \textbf{(E)}\ 64$

 \textbf{Answer: }C



\textbf{Solution 1}

By the remainder theorem, $4x^2-6x+m$ is divisible by $x-3$ if and only if $4\cdot 3^2 - 6\cdot 3 + m = 18 + m = 0$.  So $m = -18$, which is a divisor of only $\boxed{\textbf{(C) }36}$.

 -j314andrews

 



\textbf{Solution 2}

Since the given expression is a quadratic, the factored form would be $(x-3)(4x+y)$, where $y$ is a value such that $-12x+yx=-6x$ and $-3(y)=m$. The only number that fits the first equation is $y=6$, so $m=-18$. The only choice that is a multiple of 18 is $\boxed{\textbf{(C) }36}$.

 


	\item (\textit{AHSME 1953}) The base of an isosceles triangle is $6$ inches and one of the equal sides is $12$ inches. 
The radius of the circle through the vertices of the triangle is: 


 $\textbf{(A)}\ \frac{7\sqrt{15}}{5} \qquad \textbf{(B)}\ 4\sqrt{3} \qquad \textbf{(C)}\ 3\sqrt{5} \qquad \textbf{(D)}\ 6\sqrt{3}\qquad \textbf{(E)}\ \text{none of these}$

 \textbf{Answer: }E



\textbf{Solution}


                    \begin{figure}[H]
                        \centering
                        \includegraphics[width=0.5\textwidth]{Solution/AHSME_I_1953_37_Solution_Figure_1.png}
                    \end{figure}
                    Let $\triangle ABC$ be an isosceles triangle with $AB=AC=12,$ and $BC = 6$. Draw altitude $\overline{AD}$. Since $A$ is the apex of the triangle, the altitude $\overline{AD}$ is also a median of the triangle. Therefore, $BD=CD=3$. Using the \href{https://artofproblemsolving.com/wiki/index.php?title=Pythagorean\_Theorem}{Pythagorean Theorem}, $AD=\sqrt{12^2-3^2}=3\sqrt{15}$. The area of $\triangle ABC$ is $\frac12\cdot 6\cdot 3\sqrt{15}=9\sqrt{15}$.

 If $a,b,c$ are the sides of a triangle, and $A$ is its area, the circumradius of that triangle is $\frac{abc}{4A}$. Using this formula, we find the circumradius of $\triangle ABC$ to be $\frac{6\cdot 12\cdot 12}{4\cdot 9\sqrt{15}}=\frac{2\cdot 12}{\sqrt{15}}=\frac{24\sqrt{15}}{15}=\frac{8\sqrt{15}}{5}$. The answer is $\boxed{\textbf{(E)}\ \text{none of these}}$



	\item (\textit{AHSME 1953}) If $f(a)=a-2$ and $F(a,b)=b^2+a$, then $F(3,f(4))$ is: 


 $\textbf{(A)}\ a^2-4a+7 \qquad \textbf{(B)}\ 28 \qquad \textbf{(C)}\ 7 \qquad \textbf{(D)}\ 8 \qquad \textbf{(E)}\ 11$

 \textbf{Answer: }C



\textbf{Solution}

We find $f(4)=(4)-2=2$, so $F(3,f(4))=F(3,2)=(2)^2+3=\boxed{\textbf{(C) }7}$.

 


	\item (\textit{AHSME 1953}) The product, $\log_a b \cdot \log_b a$ is equal to: 


 $\textbf{(A)}\ 1 \qquad \textbf{(B)}\ a \qquad \textbf{(C)}\ b \qquad \textbf{(D)}\ ab \qquad \textbf{(E)}\ \text{none of these}$

 \textbf{Answer: }A



\textbf{Solution 1}

Let $x = \log_a b$ and $y = \log_b a$.  Then $a^x = b$ and $b^y = a$.  So $b^y = \left(a^x\right)^y = a^{xy} = a$ and $xy = \boxed{\textbf{(A) }1}$.

 



\textbf{Solution 2}

By the change-of-base formula, $\log_a b \cdot \log_b a = \frac{\log b}{\log a} \cdot \frac{\log a}{\log b} = \boxed{\textbf{(A) }1}$.

 


	\item (\textit{AHSME 1953}) The negation of the statement "all men are honest," is: 


 $\textbf{(A)}\ \text{no men are honest} \qquad \textbf{(B)}\ \text{all men are dishonest} \\  \textbf{(C)}\ \text{some men are dishonest}\qquad \textbf{(D)}\ \text{no men are dishonest}\\ \textbf{(E)}\ \text{some men are honest}$

 \textbf{Answer: }C



\textbf{Solution}

A statement and its negation cannot both be true or both be false at the same time.  Only $\boxed{\textbf{(C)}\ \text{some men are dishonest}}$ qualifies as such.

 Choice A is incorrect, as if some but not all men are honest, both this statement and the original statement are false.

 Choice B is incorrect, as if some but not all men are honest, both this statement and the original statement are false.

 Choice D is incorrect, as it is logically equivalent to the original statement.  That is, this statement and the original statement must be either both false or both true.

 Choice E is incorrect, as if no men are honest, both this and the original statement are false.

 


	\item (\textit{AHSME 1953}) A girls' camp is located $300$ rods from a straight road. On this road, a boys' camp is located $500$ rods from the girls' camp. 
It is desired to build a canteen on the road which shall be exactly the same distance from each camp. 
The distance of the canteen from each of the camps is: 


 $\textbf{(A)}\ 400\text{ rods} \qquad \textbf{(B)}\ 250\text{ rods} \qquad \textbf{(C)}\ 87.5\text{ rods} \qquad \textbf{(D)}\ 200\text{ rods}\\ \textbf{(E)}\ \text{none of these}$

 \textbf{Answer: }E



\textbf{Solution 1}


                    \begin{figure}[H]
                        \centering
                        \includegraphics[width=0.5\textwidth]{Solution/AHSME_I_1953_41_Solution_Figure_1.png}
                    \end{figure}
                    Let $r$ be the straight road, $G$ be the girls' camp, $B$ be the boys' camp, and $C$ be the water canteen. $\overline{AG}$ is the perpendicular from $G$ to $r$. Suppose each rod is one unit long. $AG=300$ and $BG=500$. Since $\angle GAB$ is a right angle, $\triangle GAB$ is a $3-4-5$ right triangle, so $AB=400$.

 Let $x$ be the distance from the canteen to the girls' and boys' camps. We have $BC=CG=x$ and $AC=400-x$. Using the \href{https://artofproblemsolving.com/wiki/index.php?title=Pythagorean\_Theorem}{Pythagorean Theorem} on $\triangle GAC$, we have
 \[300^2+(400-x)^2=x^2\] Simplifying the left side gives
 \[90000+160000-800x+x^2=x^2\]Subtracting $x^2$ from both sides and rearranging gives
 \[800x=250000\]Therefore, $x=\frac{250000}{800}=312.5$. The answer is $\boxed{\textbf{(E)}\ \text{none of these}}$.

 



\textbf{Solution 2}


                    \begin{figure}[H]
                        \centering
                        \includegraphics[width=0.5\textwidth]{Solution/AHSME_I_1953_41_Solution_Figure_2.png}
                    \end{figure}
                    

 Let $r$ be the road, $G$ be the girls' camp, $B$ be the boys' camp, $C$ be the water canteen, $A$ be the foot of the perpendicular from $G$ to $r$, and $M$ be the midpoint of $\overline{BG}$.  Since $BG = 500$ and $AG = 300$, by the Pythagorean theorem $AB = \sqrt{500^2-300^2}=400$.  Since $M$ is the midpoint of $\overline{BG}$, $MB = 250$.  Also, since $\triangle GCB$ is isosceles, $\angle CMB$ is a right angle.  So $\triangle GAB \sim \triangle CMB$ and thus $\frac{MB}{BC} = \frac{AB}{BG}$, that is, $\frac{250}{BC} = \frac{400}{500}$. Solving yields $BC = 312.5$, which is $\boxed{\textbf{(E)}\ \text{none of these}}$.

 


	\item (\textit{AHSME 1953}) The centers of two circles are $41$ inches apart. The smaller circle has a radius of $4$ inches and the larger one has a radius of $5$ inches. 
The length of the common internal tangent is: 


 $\textbf{(A)}\ 41\text{ inches} \qquad \textbf{(B)}\ 39\text{ inches} \qquad \textbf{(C)}\ 39.8\text{ inches} \qquad \textbf{(D)}\ 40.1\text{ inches}\\ \textbf{(E)}\ 40\text{ inches}$

 \textbf{Answer: }E



\textbf{Solution}


                    \begin{figure}[H]
                        \centering
                        \includegraphics[width=0.5\textwidth]{Solution/AHSME_I_1953_42_Solution_Figure_1.png}
                    \end{figure}
                    

 Let $A$ be the center of the circle with radius $5$, and $B$ be the center of the circle with radius $4$. Let $\overline{CD}$ be the common internal tangent of circle $A$ and circle $B$. Extend $\overline{BD}$ past $D$ to point $E$ such that $\overline{BE}\perp\overline{AE}$. Since $\overline{AC}\perp\overline{CD}$ and $\overline{BD}\perp\overline{CD}$, $ACDE$ is a rectangle. Therefore, $AC=DE$ and $CD=AE$.

 Since the centers of the two circles are $41$ inches apart, $AB=41$. Also, $BE=4+5=9$. Using the \href{https://artofproblemsolving.com/wiki/index.php?title=Pythagorean\_Theorem}{Pythagorean Theorem} on right triangle $ABE$, $CD=AE=\sqrt{41^2-9^2}=\sqrt{1600}=40$. The length of the common internal tangent is $\boxed{\textbf{(E) } 40\text{ inches}}$



	\item (\textit{AHSME 1953}) If the price of an article is increased by percent $p$, then the decrease in percent of sales must not exceed $d$ in order to yield the same income. 
The value of $d$ is: 


 $\textbf{(A)}\ \frac{1}{1+p} \qquad \textbf{(B)}\ \frac{1}{1-p} \qquad \textbf{(C)}\ \frac{p}{1+p} \qquad \textbf{(D)}\ \frac{p}{p-1}\qquad \textbf{(E)}\ \frac{1-p}{1+p}$

 \textbf{Answer: }C



\textbf{Solution}

It turns out that none of the listed solutions are the true solution for the problem as written, so we'll edit it slightly by replacing "percent" with "proportion". To solve the modified problem, note that the price of the article is $1+p$ times what it was originally, so that demand must be at least $\frac{1}{1+p}$ times what it was originally to yield at least the same income. Since $1-d = \frac{1}{1+p}$, this implies that $d = \frac{p}{1+p}$. Thusly, our answer is $\boxed{\text{(C)}}$, and we are done.

 


	\item (\textit{AHSME 1953}) In solving a problem that reduces to a quadratic equation one student makes a mistake only in the constant term of the equation and 
obtains $8$ and $2$ for the roots. Another student makes a mistake only in the coefficient of the first degree term and 
find $-9$ and $-1$ for the roots. The correct equation was: 


 $\textbf{(A)}\ x^2-10x+9=0 \qquad \textbf{(B)}\ x^2+10x+9=0 \qquad \textbf{(C)}\ x^2-10x+16=0\\  \textbf{(D)}\ x^2-8x-9=0\qquad \textbf{(E)}\ \text{none of these}$

 \textbf{Answer: }A



\textbf{Solution}

Let $x^2+bx+c=0$ represent the correct equation. Since the coefficient of the $x^2$ term is $1$, the sum of the roots is $-b$ and the product of the roots is $c$.

 The first student has the correct first-degree term, so the sum of the roots is $8+2=10$ and $b=-10$. The second student has the correct constant term, so the product of the roots is $(-9)(-1) = 9$ and $c=9$. Therefore, the correct equation is $\boxed{\textbf{(A) } x^2-10x+9=0}$.

 


	\item (\textit{AHSME 1953}) The lengths of two line segments are $a$ units and $b$ units respectively. Then the correct relation between them is: 


 $\textbf{(A)}\ \frac{a+b}{2} > \sqrt{ab} \qquad \textbf{(B)}\ \frac{a+b}{2} < \sqrt{ab} \qquad \textbf{(C)}\ \frac{a+b}{2}=\sqrt{ab}\\ \textbf{(D)}\ \frac{a+b}{2}\leq\sqrt{ab}\qquad \textbf{(E)}\ \frac{a+b}{2}\geq\sqrt{ab}$

 \textbf{Answer: }E



\textbf{Solution}

Since $a$ and $b$ are positive, the \href{https://artofproblemsolving.com/wiki/index.php?title=AM-GM\_Inequality}{AM-GM Inequality}, that is, $\boxed{\textbf{(E)}\ \frac{a+b}{2}\geq\sqrt{ab}}$, holds.

 Choice A is false if $a = b$.  Choice B is always false.  Choices C and D are false if $a \neq b$.

 


	\item (\textit{AHSME 1953}) Instead of walking along two adjacent sides of a rectangular field, a boy took a shortcut along the diagonal of the field and 
saved a distance equal to $\frac{1}{2}$ the longer side. The ratio of the shorter side of the rectangle to the longer side was: 


 $\textbf{(A)}\ \frac{1}{2} \qquad \textbf{(B)}\ \frac{2}{3} \qquad \textbf{(C)}\ \frac{1}{4} \qquad \textbf{(D)}\ \frac{3}{4}\qquad \textbf{(E)}\ \frac{2}{5}$

 \textbf{Answer: }D



\textbf{Solution}

Let $x<y$ be the sides of the rectangle. The length of the diagonal is $\sqrt{x^2+y^2}$, and the length of the two adjacent sides is $x+y$. Then the distance the boy saves is $x+y-\sqrt{x^2+y^2}$. Setting this equal to $\frac12y$, we have
 \[x+y-\sqrt{x^2+y^2}=\frac12y\]
 \[x+\frac12y=\sqrt{x^2+y^2}\]
 \[x^2+xy+\frac14y^2=x^2+y^2\]
 \[xy=\frac34y^2\]
 \[\frac xy=\frac34,\]so the answer is $\boxed{\textbf{(D)}\ \frac{3}{4}}$.

 


	\item (\textit{AHSME 1953}) If $x>0$, then the correct relationship is: 


 $\textbf{(A)}\ \log (1+x) = \frac{x}{1+x} \qquad \textbf{(B)}\ \log (1+x) < \frac{x}{1+x} \\  \textbf{(C)}\ \log(1+x) > x\qquad \textbf{(D)}\ \log (1+x) < x\qquad \textbf{(E)}\ \text{none of these}$

 \textbf{Answer: }D



\textbf{Solution 1 (Calculus)}

If $x=9$, then choices A, B, and C are incorrect.  

 Let $f(x) = x - \ln(1+x)$.  Then $f'(x) = 1 - \frac{1}{1+x} = \frac{x}{1+x}$, which must be positive if $x > 0$.  Thus $f(x)$ must be increasing while $x > 0$ and $f(x) > 0$ whenever $x > 0$.

 So $x > \ln(1+x) = \frac{\log(1+x)}{\log{e}} > \log(1+x)$, so the correct answer is $\boxed{\textbf{(D)} \log (1+x) < x}$.

 -j314andrews

 



\textbf{Solution 2 (Taylor Series)}

Recall that the Taylor series of $e^x$ is $1 + x + \frac{x^2}{2} + \frac{x^3}{6} + ...$, which must be greater than $1 + x$. So $1 + x < e^x < {10}^x$ and $\boxed{\textbf{(D)} \log (1+x) < x}$.

 -j314andrews

 


	\item (\textit{AHSME 1953}) If the larger base of an isosceles trapezoid equals a diagonal and the smaller base equals the altitude, 
then the ratio of the smaller base to the larger base is: 


 $\textbf{(A)}\ \frac{1}{2} \qquad \textbf{(B)}\ \frac{2}{3} \qquad \textbf{(C)}\ \frac{3}{4} \qquad \textbf{(D)}\ \frac{3}{5}\qquad \textbf{(E)}\ \frac{2}{5}$

 \textbf{Answer: }D



\textbf{Solution}


                    \begin{figure}[H]
                        \centering
                        \includegraphics[width=0.5\textwidth]{Solution/AHSME_I_1953_48_Solution_Figure_1.png}
                    \end{figure}
                    

 Let $a$ be the length of the smaller base of isosceles trapezoid $ABCD$, and $1$ be the length of the larger base of the trapezoid. The ratio of the smaller base to the larger base is $\frac a1=a$. Let point $E$ be the foot of the altitude from $C$ to $\overline{AD}$.

 Since the larger base of the isosceles trapezoid equals a diagonal, $AC=AD=1$. Since the smaller base equals the altitude, $BC=CE=a$. Since the trapezoid is isosceles, $DE=\frac{1-a}{2}$, so $AE = 1-\frac{1-a}{2} = \frac{a+1}{2}$. Using the \href{https://artofproblemsolving.com/wiki/index.php?title=Pythagorean\_Theorem}{Pythagorean Theorem} on right triangle $ACE$, we have
 \[a^2+\left(\frac{a+1}{2}\right)^2=1\]Multiplying both sides by $4$ gives
 \[4a^2+(a+1)^2=4\]Expanding the squared binomial and rearranging gives
 \[5a^2+2a-3=0\]This can be factored into $(5a-3)(a+1)=0$. Since a must be positive, $5a+3=0$, so $a=\frac 35$. The ratio of the smaller base to the larger base is $\boxed{\textbf{(D) } \frac 35}$.

 


	\item (\textit{AHSME 1953}) The coordinates of $A,B$ and $C$ are $(5,5),(2,1)$ and $(0,k)$ respectively. 
The value of $k$ that makes $\overline{AC}+\overline{BC}$ as small as possible is: 


 $\textbf{(A)}\ 3 \qquad \textbf{(B)}\ 4\frac{1}{2} \qquad \textbf{(C)}\ 3\frac{6}{7} \qquad \textbf{(D)}\ 4\frac{5}{6}\qquad \textbf{(E)}\ 2\frac{1}{7}$

 \textbf{Answer: }E



\textbf{Solution}

Let $A'$ be the reflection of $A$ over the $y$-axis. Then $A'$ has coordinates $(-5,5)$ and $A'C = AC$.  So $A'C+BC=AC+BC$. By the triangle inequality, $A'C+BC \geq A'B$, with equality happening when $C$ lies on segment ${A'B}$.  Therefore, $k$ is the $y$-intercept of line $A'B$.  

 The slope of line $A'B$ is $\frac{1-5}{2-(-5)}=-\frac {4}{7}$. In point-slope form, the equation of this line is $y-5=-\frac {4}{7}(x+5)$, which becomes $y = -\frac{4}{7}x + \frac{15}{7}$ when converted to slope-intercept form.  So $k = \frac{15}{7} = \boxed{\textbf{(E) } 2\frac17}.$



	\item (\textit{AHSME 1953}) One of the sides of a triangle is divided into segments of $6$ and $8$ units by the point of tangency of the inscribed circle. 
If the radius of the circle is $4$, then the length of the shortest side of the triangle is: 


 $\textbf{(A)}\ 12\text{ units} \qquad \textbf{(B)}\ 13\text{ units} \qquad \textbf{(C)}\ 14\text{ units} \qquad \textbf{(D)}\ 15\text{ units}\qquad \textbf{(E)}\ 16\text{ units}$

 \textbf{Answer: }B



\textbf{Solution 1}

Let the triangle have side lengths $14, 6+x,$ and $8+x$. The area of this triangle can be computed two ways. We have $A = rs$, and $A = \sqrt{s(s-a)(s-b)(s-c)}$, where $s = 14+x$ is the semiperimeter. Therefore, $4(14+x)=\sqrt{(14+x)(x)(8)(6)}$. Solving gives $x = 7$ as the only valid solution. This triangle has sides $13,14$ and $15$, so the shortest side is $\boxed{\textbf{(B) \ } 13 \mathrm{\ units}}$.

 



\textbf{Solution 2}

Label the tangent points on $BC, CA, AB$ as $D, E, F$ respectively. Let $AF=AE=6$, $BF=BD=8$, and $CE=CD=x.$ The problem is a matter of solving for $x$. To this, we use the fact that if $A,B,C$ are the angles of a triangle, then $\tan{\frac{A}{2}}\tan{\frac{B}{2}}+\tan{\frac{B}{2}}\tan{\frac{C}{2}}+\tan{\frac{C}{2}}\tan{\frac{A}{2}} = 1.$ We know that $\tan{\frac{A}{2}} = \frac{2}{3}$, $\tan{\frac{B}{2}} = \frac{1}{2}$, and $\tan{\frac{C}{2}} = \frac{4}{x},$ so we have the equation $\frac{1}{2}\cdot \frac{2}{3} + \frac{1}{2}\cdot \frac{4}{x} + \frac{4}{x}\cdot \frac{2}{3} = 1.$ Solving this equation yields $x=7$, so the shortest side has length $\boxed{\textbf{(B) \ } 13 \mathrm{\ units}}$.

 


\end{enumerate}